\section{课程主线：AGI 不是一个形状，而是一套可承担责任的系统}

“The Advent of AGI” 很容易被读成关于时间表、能力上限或超级智能的预测。Div Garg 的课堂主线更工程化：即使模型已经能聊天、推理和调用工具，\textbf{把智能变成日常可用的行动系统}仍需要定义任务、权限、记忆、训练、评测、通信与失败恢复。AGI 的 form factor 尚未确定，因此不能先假设它一定是更大的聊天窗口；应先问用户把什么责任交给系统、系统凭什么被信任、失败时谁能接管。

\lecturefigure{slide-01-title-advent-agi.jpg}{课程标题：The Advent of AGI}{00:00:15}

\lecturefigure{slide-02-what-does-agi-look-like.jpg}{开场问题：AGI 到底以什么形态进入生活}{00:01:50}

\teachervoice{00:01:18--00:02:10，讲者说聊天与推理能力已经迅速提高，但“AGI 看起来像什么”仍没有共同答案。课堂提示：能力水平与产品形态是两个问题；后者决定系统如何进入社会、如何被验证以及谁承担风险。}

\begin{importantbox}{本讲的操作性定义}
本讲不把 AGI 定义为某个参数规模或 benchmark 阈值，而把它视为一个逐渐扩大的\term{责任边界}：系统能感知环境、形成计划、使用工具、持续记忆并代表用户行动。能力每扩大一步，都必须同时扩大评测、权限、观测和恢复机制。
\end{importantbox}

\subsection{Agent 的四层契约：Memory、Tools、Planning、Action}

第一张系统图来自 Lilian Weng 对 LLM-powered agents 的经典拆分。读图时不要只记四个名词，而要看信息怎样流动：短期上下文和长期记忆给出状态；工具把语言连接到外部世界；规划产生可执行子目标；动作改变环境并带回新观察。任何一层缺失，系统都会退化成无法持续闭环的单次生成器。

\lecturefigure{slide-03-agent-architecture-components.jpg}{Agent 架构：记忆、工具、规划与动作共同形成闭环}{00:03:53}

可把一个工具型 agent 写成部分可观测决策过程（Partially Observable Markov Decision Process, POMDP）：
\[
\mathcal{M}=(\mathcal{S},\mathcal{A},P,R,\Omega,O,\gamma).
\]
其中 $\mathcal{S}$ 是真实环境状态，$\mathcal{A}$ 是动作集合，$P$ 是状态转移，$R$ 是奖励，$\Omega$ 是可见观察，$O$ 描述状态怎样生成观察，$\gamma$ 是折扣因子。模型通常只能看到网页、消息、工具返回和压缩记忆构成的 $o_t\in\Omega$，而不是完整的 $s_t\in\mathcal{S}$；因此“看起来合理”的下一步不一定是全局正确动作。

\begin{knowledgebox}{首次使用：Tool、Memory 与 Planning 分别解决什么}
\textbf{Tool} 是模型可调用的外部能力，例如浏览器、数据库、日历或代码执行器；\textbf{memory} 是跨步骤保存和检索状态的机制，不等于把所有历史塞进 prompt；\textbf{planning} 是把目标分解为子目标、检查进展并在失败后改路。三者的接口必须显式，否则错误无法定位到“模型不会”“工具失败”还是“状态丢失”。
\end{knowledgebox}

\begin{warningbox}{把 CoT 当作计划的危险}
Chain-of-Thought（CoT）是语言化推理轨迹，不自动等于可执行计划。真实计划还需要动作前置条件、权限、预算、终止条件、幂等性和回滚。一个语言上连贯的步骤序列，可能在第三步就已经违反环境约束。
\end{warningbox}

\subsection{Everyday AGI：愿景必须立刻翻译成研究议程}

讲者用 Everyday AGI 和 DMV 驾驶考试演示表达一种产品愿景：AI 不只是回答问题，而是能够替用户完成现实流程。画面本身只能证明当时存在一个课堂 demo，不能证明所有州、所有考试或所有网站均可稳定完成。更重要的是，下一张 slide 立即把愿景拆成 evaluation、training 与 communication 三条研究线，说明产品承诺必须被转成可检验的系统问题。

\lecturefigure{slide-04-everyday-agi.jpg}{Everyday AGI：把智能嵌入日常流程的产品愿景}{00:04:11}

\lecturefigure{slide-05-dmv-agent-demo.jpg}{课堂展示的 DMV 在线考试代理场景}{00:05:28}

\lecturefigure{slide-06-agent-initiatives.jpg}{从愿景到研究议程：评测、训练与 Agent 通信}{00:05:44}

这三张图共同展示了“愿景--能力--基础设施”的递进。Everyday AGI 给出用户可感知的目标，DMV 画面给出一个具体交互样本，initiatives 则列出让样本变成可重复系统所需的研究工作。若只保留第一层，课程会退化成未来叙事；若只保留第三层，又会失去为什么需要这些技术。产品团队应把每个高影响 demo 写成任务规范，再反推需要哪些 deterministic environments、训练轨迹和跨 agent 协议。

\teachervoice{00:04:11--00:07:10，讲者先展示“日常 AGI”愿景，再列出 REAL Evals、AgentQ 与 agent communication。讲义提醒：demo 是提出问题的入口，不是可靠性证书；真正的主线是怎样评、怎样学、怎样协作。}

\begin{importantbox}{从 Demo 到系统验收的五个问题}
\begin{enumerate}
  \item 任务是否有可判定的完成条件？
  \item 环境变化后能否复现同一测试？
  \item 中间动作是否全部可追踪？
  \item 失败能否恢复而不是重头盲试？
  \item 高风险动作是否需要用户确认或人工接管？
\end{enumerate}
\end{importantbox}

\subsection{本章小结}

本章把“AGI 到来”压缩为一个系统合同：模型必须在部分可观测环境中，用记忆、工具、规划与行动完成闭环；而任何产品愿景都要被立即翻译成评测、训练、通信和责任边界。下一章进一步回答：为什么要采用“human-like”计算机交互，以及自主程度怎样分级。

